\documentclass[12pt, twoside]{article}
%\documentclass[12pt, twoside]{article}
\usepackage[letterpaper, margin=1in, headsep=0.2in]{geometry}
\setlength{\headheight}{0.6in}
%\usepackage[english]{babel}
\usepackage[utf8]{inputenc}
\usepackage{microtype}
\usepackage{amsmath}
\usepackage{amssymb}
%\usepackage{amsfonts}
\usepackage[nomessages]{fp} %\FPeval{\var-name}{2*sin(pi/6)}
\usepackage{siunitx} %units in math. eg 20\milli\meter
\usepackage{yhmath} % for arcs, overparenth command
\usepackage{tikz} %graphics
\usetikzlibrary{quotes, angles, arrows, arrows.meta}
\usepackage{pgfplots}
\usepgfplotslibrary{fillbetween}
\pgfplotsset{compat=1.16}
\usepackage{graphicx} %consider setting \graphicspath{{images/}}
\usepackage{parskip} %no paragraph indent
\usepackage{enumitem}
\usepackage{multicol}
\usepackage{venndiagram}

\usepackage{fancyhdr}
\pagestyle{fancy}
\fancyhf{}
\renewcommand{\headrulewidth}{0pt} % disable the underline of the header
\raggedbottom
\hfuzz=2mm %suppresses overfull box warnings

\usepackage{hyperref}
\setlength{\parskip}{0.3\baselineskip}

\fancyhead[LE]{\thepage}
\fancyhead[RO]{Markscheme}
\fancyhead[LO]{La Scuola d'Italia / Dr.\ Huson / IB Math 12 \\* 5 October 2026}

\begin{document}
\subsubsection*{1.7 Classwork: Cumulative frequency graphs --- Markscheme}

\emph{Calculator allowed. Marks are not printed on the student sheet; they are
recorded here. Total: 17 marks.}

\emph{\textbf{Throughout:} a value read from a graph is accepted within the
range given. Lines drawn on the graph that show the reading count as working
for an M1. Where a later part uses an earlier reading, follow through (FT)
from the student's value.}

\begin{enumerate}[itemsep=0.3cm]

%--- Problem 1: Read median, quartiles, IQR, a count and a percentile from a CF curve; outlier test [9 marks] ---
\item[1.] [Maximum mark: 9]

$100$ students; curve passes through $(30, 15)$, $(40, 25)$, $(60, 50)$,
$(80, 75)$, $(100, 90)$.

\textbf{(a)} Median at cumulative frequency $50$: $\;60$ minutes
\hfill \textbf{A1}

\emph{Accept $59$ to $61$.}

\textbf{(b)} $Q_{1} = 40$ and $Q_{3} = 80$ \hfill \textbf{A1}

$\text{IQR} = 80 - 40 = 40$ minutes \hfill \textbf{A1}

\emph{Accept $Q_{1}$ in $39$ to $41$ and $Q_{3}$ in $79$ to $81$. The second
A1 is FT from the student's quartiles. Reading the quartiles at cumulative
frequencies other than $25$ and $75$ earns A0A0.}

\textbf{(c)} $15$ students spent less than $30$ minutes \hfill
\textbf{(M1)}

$100 - 15 = 85$ students \hfill \textbf{A1}

\emph{Accept $84$ to $86$. The M1 is for reading the curve at $t = 30$,
seen or implied by a line on the graph. ``$15$'' as the final answer earns
M1A0: it is the number below $30$ minutes, not at least $30$.}

\textbf{(d)} $90\%$ of $100 = 90$ students, so read the curve at cumulative
frequency $90$ \hfill \textbf{(M1)}

$100$ minutes \hfill \textbf{A1}

\emph{Accept $99$ to $101$. Reading at cumulative frequency $10$ (the
$10^{\text{th}}$ percentile, about $24$ minutes) earns M0A0.}

\textbf{(e)} $Q_{1} - 1.5 \times \text{IQR} = 40 - 60 = -20$ \quad and \quad
$Q_{3} + 1.5 \times \text{IQR} = 80 + 60 = 140$ \hfill \textbf{M1}

$5 > -20$ and $118 < 140$, so neither the shortest nor the longest time is
an outlier \hfill \textbf{R1}

no outliers \hfill \textbf{AG}

\emph{M1 for at least one boundary correctly calculated (FT from (b)). The R1
needs both comparisons, with the minimum and the maximum. Checking only the
longest time earns M1R0.}

\newpage

%--- Problem 2: Read median and a count from a CF curve; complete a grouped frequency table; estimate the mean [8 marks] ---
\item[2.] [Maximum mark: 8]

$120$ apples; curve passes through $(120, 10)$, $(140, 35)$, $(160, 75)$,
$(170, 90)$, $(180, 105)$, $(200, 120)$.

\textbf{(a)} Median at cumulative frequency $60$: $\;152.5$ grams
\hfill \textbf{A1}

\emph{Accept $151$ to $154$. Reading at cumulative frequency $60.5$ is also
accepted.}

\textbf{(b)} $90$ apples have a mass of $170$ g or less \hfill
\textbf{(M1)}

$120 - 90 = 30$ large apples \hfill \textbf{A1}

\emph{Accept $29$ to $31$. ``$90$'' as the final answer earns M1A0.}

\textbf{(c)} $p = 75 - 35 = 40$ \hfill \textbf{A1}

$q = 105 - 75 = 30$ \hfill \textbf{A1}

\emph{Check: $10 + 25 + 40 + 30 + 15 = 120$. Cumulative frequencies read
as the frequencies ($p = 75$, $q = 105$) earn A0A0.}

\textbf{(d)} Mid-interval values $110, 130, 150, 170, 190$ with frequencies
$10, 25, 40, 30, 15$ \hfill \textbf{(M1)}

$\bar{x} = \dfrac{1100 + 3250 + 6000 + 5100 + 2850}{120} =
\dfrac{18\,300}{120} = 152.5$ grams \hfill \textbf{A1}

\emph{\textbf{GDC:} mid-interval values in L1, frequencies in L2, One-Variable
Statistics with frequency list L2: $\bar{x} = 152.5$.}

\emph{FT from the student's $p$ and $q$. Using the class boundaries ($120,
140, \ldots$) instead of the mid-interval values earns M0A0. The mean of the
five mid-interval values alone ($150$) earns M0A0.}

\textbf{(e)} The actual masses are not known; each apple is given the
mid-interval value of its class. \hfill \textbf{R1}

\emph{Accept ``the data are grouped'', ``mid-interval values are used''.
``The graph is not exact'' alone earns R0: the mean is calculated from the
table, and it would be an estimate even with exact frequencies.}

\end{enumerate}
\end{document}
